\begin{proposition}\label{prop:nonic33-finite-classes}
Put $X=X_1(33)$. Every rational effective divisor $D$ of degree $9$
on $X$ is either one of the following known divisors, or is linearly
equivalent over $\mathbb Q$ to one of the ten signed cuspidal divisors
$C_i$ in \texttt{short\_cusp\_classes33.json}:
\begin{enumerate}
\item a rational effective cuspidal divisor;
\item a known septic divisor plus a rational effective cuspidal divisor
of degree $2$;
\item a known octic divisor plus a rational cusp;
\item a known closed nonic divisor in the diamond orbit of the
explicit nonic $33$-torsion example.
\end{enumerate}
Each $C_i$ has degree $9$ and is an integer combination of global
rational closed cusps of degrees $1$ and $2$.
\end{proposition}
\begin{proof}
The rational gonality of $X$ is $10$ by \cite[Table~1]{DvH}.
By \cite[Theorem~3.1(2), Theorem~4.13 and Table~2]{DEHMZ},
\[
 J_1(33)(\mathbb Q)\simeq C_5\times C_{42373650}
\]
is generated by rational cuspidal divisors. This published global
upper bound is independent of the reduction computation below.
An unconditional finite
class-group computation gives
\[
 J_1(33)(\mathbb F_2)\simeq C_{15}\times C_{42373650}.
\]
The reductions of the global cuspidal generators generate a subgroup
$H$ of invariants $[5,42373650]$. Thus reduction on the full rational
Jacobian is injective, including on its $2$-primary subgroup, and
$H$ is its complete image.

The rational cusp-orbit degrees are $1^{10},2^5,10^2$, and their
reductions have closed-place degrees $1^{10},2^5,5^2,10$.
The signatures
$\bigl(v_c(b^4/\Delta),v_c(b)\bigr)$ distinguish the $17$ global
orbits. The two functions are integral modular units away from $33$;
their divisors have no vertical components there. Hence these
signatures identify the same global closed cusp divisors in every
good characteristic. In particular no individual finite part of a
split nonrational global cusp orbit is used as a global divisor.

The noncuspidal place counts through degree $9$ over $\mathbb F_2$
are
\[
 (0,0,0,0,2,20,30,25,50).
\]
Since every noncuspidal component has degree at least $5$, an
effective divisor of degree $9$ has at most one such component.
All multiplicities in its cusp complement are retained.
The complete class tests are
\[
\begin{array}{c|r|r|r}
\text{noncuspidal degree}&\text{tests}&\text{class in }H&
 \text{known classes}\\\hline
0&146260&144250&144250\\
5&2010&0&0\\
6&5400&0&0\\
7&1800&600&600\\
8&250&50&50\\
9&50&20&10.
\end{array}
\]
The $144250$ first-row classes are precisely the rational effective
cuspidal classes. The next two nonempty known sets consist of the
$10$ septic diamond translates with each of the $60$ rational cusp
complements of degree $2$, and the $5$ octic translates with each of
the $10$ rational cusps. The remaining ten known classes are the
nonic diamond orbit. All these global witnesses are reconstructed
and their point orders checked in the computation.

The other ten classes are represented by the stated $C_i$.
Their positive and negative degrees are
\[
 (15,6),(15,6),(17,8),(15,6),(13,4),
 (15,6),(14,5),(13,4),(14,5),(14,5).
\]
The integer combinations use the complete global cusp blocks.
Their degrees and finite class identities are independently checked.
For any possible reduction, injectivity of the rational Jacobian
identifies the global class with the corresponding known divisor or
with $C_i$. In the known effective cases, gonality $10$ then
identifies the effective divisors themselves. This proves the claim.

The independent replay verifies all $155770$ class tests, all $127$
finite noncuspidal point coordinates (which are nonsingular affine
points), the full rational subgroup, every complement, and all ten
cuspidal class identities.
\end{proof}

\begin{lemma}\label{lem:nonic33-second-reduction}
If $h^0(X_{\mathbb F_q},\overline C_i)=0$ at any prime $q\nmid33$,
then no rational effective divisor of degree $9$ on $X$ is linearly
equivalent to $C_i$.
\end{lemma}
\begin{proof}
The global cusps extend as closed divisors on the smooth proper model
at $q$. Any effective divisor in the characteristic-zero class of
$C_i$ would specialize to an effective divisor in the class of
$\overline C_i$, and would give a nonzero section of its line bundle.
This contradicts the asserted vanishing.
\end{proof}

\begin{theorem}\label{thm:nonic33-complete}
Every rational effective divisor of degree $9$ on $X_1(33)$ supported
away from the cusps is one of the ten known nonic divisors.
If an elliptic curve $E$ over a degree-$9$ number field $K$ has a point
of order $33$, then $E(K)_{\mathrm{tors}}\simeq C_{33}$.
In particular, $C_{66}\notin\Phi(9)$.
\end{theorem}
\begin{proof}
The ten residual classes of Proposition~\ref{prop:nonic33-finite-classes}
are removed by Lemma~\ref{lem:nonic33-second-reduction} at $q=5$.
Here is an exact certificate for the required vanishing.
Put $A=\sum_c\max_i\{0,v_c(C_i)\}\,c$, taking maxima over the ten
explicit global cuspidal combinations. Then $\deg A=48$ and
$B_i=A-C_i$ is effective of degree $39$ for every $i$.
Since $g(X_1(33))=21$, Riemann--Roch gives
\[
 h^0(A)=48+1-21=28.
\]
The certificate provides $28$ functions with poles bounded by $A$.
Their coefficients in the degree-$12$ function-field basis are
linearly independent over $\mathbb F_5$, so these functions form the
whole space $L(A)$.

At each place in the support of the $B_i$, expand a basis function
$f$ after multiplication by a local uniformizer to the power $v_c(A)$.
Its first $v_c(B_i)$ coefficients give the restriction to $B_i$.
Expanding residue-field coefficients over $\mathbb F_5$ gives a
$28\times39$ matrix of rank $28$ for every $i$.
Consequently $L(A-B_i)=L(C_i)=0$ in all ten cases.
A fresh verifier checks every basis pole bound and every needed
local vanishing kernel without computing a Riemann--Roch space or
a class group. A separate Sage computation verifies basis
independence and all ten matrix ranks.

The only remaining effective divisors in the finite classification
are the known cuspidal divisors, the septic-plus-cusp and
octic-plus-cusp divisors, and the ten known nonic divisors.
The first three types contain cusps, proving the first assertion.

For completeness, the nonic point has coordinates $x=a$ and
\[
 y=\frac{113a^8-638a^7-71a^6+3534a^5+544a^4
 -5545a^3-1342a^2+2083a+197}{543},
\]
where
\[
 a^9-5a^8-3a^7+25a^6+20a^5-19a^4-15a^3+4a^2+2a-1=0.
\]
This polynomial is irreducible. Its discriminant is a $5$-adic unit,
and the quadratic factor $a^2+2a-1$ modulo $5$ gives good reduction
of the associated Tate curve over $\mathbb F_{25}$.
Direct enumeration gives exactly $33$ points on this reduction.
Unramified odd-prime reduction injects torsion, and the global point
has exact order $33$. Thus the full torsion is $C_{33}$, as it is for
every diamond translate.

A point of order $33$ over $K$ gives a noncuspidal rational trace
divisor of degree $9$, allowing multiplicities if its field of
moduli is smaller. The classification forces this trace to be a
known closed nonic divisor, so its field of definition is all of $K$.
The fine moduli interpretation identifies the elliptic curve over
that field with the known curve. Its full torsion is therefore
$C_{33}$. A point of order $66$ would also give a point of order $33$,
which proves the final assertion.
\end{proof}
